\begin{table}[H]
\begin{threeparttable}
\caption{The Leapfrog Dynamic: Financial Exclusion (2014) and Fintech Adoption (2022)}
\label{atab:leapfrog}
\small\begin{tabular}{lccc}
\toprule
Geopolitical zone & Unbanked (2014, \%) & Fintech index (2022, \%) & States ($N$) \\
\midrule
3 & 38.3 & 30.3 & 7 \\
1 & 32.2 & 30.1 & 7 \\
4 & 30.7 & 25.6 & 5 \\
5 & 30.6 & 24.6 & 6 \\
6 & 21.1 & 24.4 & 6 \\
2 & 19.4 & 31.3 & 6 \\
\midrule
Pearson $r$ & \multicolumn{3}{c}{0.126} \\
\bottomrule
\end{tabular}
\begin{tablenotes}\footnotesize
\item \textit{Notes:} Zone-level means. Unbanked (2014) is the share of WBES 2014 firms without a bank account (proxy for initial financial exclusion). Fintech index (2022) is the mean state-level digital payment penetration from NLPS Round 6 (October 2022). The weak, near-zero correlation ($r=0.13$) between 2014 exclusion and 2022 fintech adoption reflects the leapfrog dynamic: zones that were most financially excluded in 2014 had the strongest incentive to adopt mobile money between 2014 and 2022, invalidating road distance to financial hubs as an instrument for fintech adoption. This explains the wrong-sign first-stage coefficient in Appendix Table~\ref{atab:iv_firststage}.
\end{tablenotes}\end{threeparttable}\end{table}
